% chapters/ch05_risk_state_prices_and_expected_utility.tex

\section{Risk, State Prices, and Expected Utility}

This chapter extends consumer theory to uncertain environments. Instead of choosing bundles across goods or across time, the consumer now chooses contingent consumption across states of the world. This allows a unified treatment of insurance, asset markets, diversification, and expected utility. The central ideas are:
\begin{enumerate}
    \item uncertain payoffs can be represented as contingent consumption bundles;
    \item state prices and asset prices are linked by no-arbitrage;
    \item diversification and risk sharing can smooth consumption across states;
    \item expected utility provides a tractable model of risk preferences.
\end{enumerate}

\subsection{Risk and Contingent Consumption}

Suppose there are $S$ mutually exclusive states of the world. State $s$ occurs with probability $\pi_s$, where
\[
\pi_s \ge 0
\qquad \text{and} \qquad
\sum_{s=1}^S \pi_s = 1.
\]

\begin{definition}{Contingent consumption bundle}{ch5-contingent-bundle}
A \emph{contingent consumption bundle} is a vector
\[
x=(x_1,x_2,\dots,x_S)\in \mathbb{R}_+^S,
\]
where $x_s$ is the amount of money or consumption received if state $s$ occurs.
\end{definition}

This is the correct representation for risky prospects. A lottery, job outcome, insurance payoff, or asset return can all be written as a contingent bundle.

\begin{definition}{State prices and the contingent-claims budget set}{ch5-state-prices}
Let
\[
p=(p_1,p_2,\dots,p_S)\in \mathbb{R}_{++}^S
\]
be the \emph{state price vector}, where $p_s$ is the price today of one unit of consumption delivered only in state $s$.

If wealth is $m>0$, the consumer chooses $x$ subject to
\[
p\cdot x = \sum_{s=1}^S p_s x_s \le m.
\]
If the consumer has endowment
\[
\omega=(\omega_1,\dots,\omega_S),
\]
then the budget constraint becomes
\[
p\cdot x \le p\cdot \omega.
\]
\end{definition}

State prices play exactly the same role here that ordinary prices played in Chapters 1--4. The difference is that the coordinates now refer to states rather than goods or dates.

\begin{examplebox}[title={Worked Example: contingent consumption from an uncertain career outcome}]
Suppose there are two states:
\begin{itemize}
    \item state 1: a high-paying job occurs with probability $0.2$;
    \item state 2: onion farming occurs with probability $0.8$.
\end{itemize}
If the corresponding incomes are
\[
10^6 \quad \text{and} \quad 10^3,
\]
then the risky prospect is represented by
\[
x=(10^6,10^3),
\qquad
\pi=(0.2,0.8).
\]
This is a contingent consumption bundle because the realized payoff depends on which state occurs.
\end{examplebox}

\subsection{Insurance and State Prices}

Insurance can be interpreted as trading contingent claims.

\begin{examplebox}[title={Worked Example: insurance as contingent consumption}]
Suppose the consumer is endowed with
\[
\omega=(8,2),
\qquad
\pi=(0.9,0.1),
\]
so state 1 is the good state and state 2 is the bad state. A unit of insurance:
\begin{itemize}
    \item costs $2$ today;
    \item pays $3$ in state 2.
\end{itemize}

Relative to the endowment, buying one unit of insurance reduces state-1 consumption by $2$ and raises state-2 consumption by $1$. If $\theta$ units are purchased, then the resulting contingent bundle is
\[
x(\theta) = (8-2\theta,\, 2+\theta).
\]

If shorting insurance is allowed, then $\theta$ may be negative. For example,
\[
\theta=-2
\]
gives
\[
x(-2)=(12,0).
\]

Now observe that
\[
x_1 + 2x_2 = (8-2\theta)+2(2+\theta)=12
\]
for every $\theta$. Since
\[
12 = 8 + 2\cdot 2 = (1,2)\cdot (8,2),
\]
this budget line is exactly the same as the contingent-claims budget with state price vector
\[
\boxed{p=(1,2).}
\]

So insurance trading is equivalent to buying state-contingent consumption at prices $(1,2)$.
\end{examplebox}

\begin{examtip}
When an insurance contract changes payoffs by a fixed amount across states, rewrite the set of attainable bundles as a linear equation. That linear equation usually reveals the implied state prices immediately.
\end{examtip}

\subsection{Assets, Portfolios, and Arbitrage}

Assets can also be viewed as contingent consumption bundles.

\begin{definition}{Assets, portfolios, and shorting}{ch5-assets}
An asset is a contingent payoff vector
\[
a^k=(a_1^k,\dots,a_S^k)\in \mathbb{R}^S.
\]
If its price is $q_k$, then holding one unit of asset $k$ costs $q_k$ today and delivers $a_s^k$ in state $s$.

A \emph{portfolio} is a vector
\[
\theta=(\theta_1,\dots,\theta_K),
\]
where $\theta_k$ is the quantity of asset $k$ held.

The portfolio costs
\[
q\cdot \theta = \sum_{k=1}^K q_k \theta_k
\]
and delivers contingent consumption
\[
\sum_{k=1}^K \theta_k a^k.
\]

If $\theta_k<0$ is allowed, then shorting is allowed.
\end{definition}

\begin{definition}{Arbitrage}{ch5-arbitrage}
A portfolio $\theta$ is an \emph{arbitrage opportunity} if:
\begin{enumerate}
    \item it costs nothing or yields money upfront:
    \[
    q\cdot \theta \le 0;
    \]
    \item its payoff is nonnegative in every state:
    \[
    \sum_{k=1}^K \theta_k a^k \ge 0;
    \]
    \item and it is strictly positive in at least one state.
\end{enumerate}
\end{definition}

Arbitrage matters because no-arbitrage is the key pricing principle of the chapter: if two portfolios have the same state-contingent payoffs, they must have the same price.

\begin{theorem}{Pricing by replication}{ch5-replication}
Suppose a new asset $\tilde a$ can be replicated by an existing portfolio $\theta$, i.e.
\[
\tilde a = \sum_{k=1}^K \theta_k a^k.
\]
If there is no arbitrage, then the price of the new asset must be
\[
\tilde q = q\cdot \theta.
\]
\end{theorem}

\begin{proof}
If $\tilde q < q\cdot \theta$, then buy the new asset and short the replicating portfolio. This gives a positive profit today and zero net payoff in every state, which is an arbitrage.

If $\tilde q > q\cdot \theta$, then short the new asset and buy the replicating portfolio. Again, the payoffs cancel state by state, leaving a sure profit today.

So no-arbitrage requires
\[
\tilde q = q\cdot \theta.
\]
\end{proof}

\begin{definition}{State-price representation}{ch5-state-representation}
If Arrow securities exist, where:
\[
e^1=(1,0,\dots,0),\ e^2=(0,1,\dots,0),\ \dots,\ e^S=(0,\dots,0,1),
\]
and their prices are
\[
p=(p_1,\dots,p_S),
\]
then any asset
\[
a=(a_1,\dots,a_S)
\]
has no-arbitrage price
\[
q(a)=\sum_{s=1}^S p_s a_s = p\cdot a.
\]
\end{definition}

So state prices are the building blocks of all asset prices.

\begin{examplebox}[title={Worked Example: extracting state prices from asset prices}]
Suppose there are two assets:
\[
a^1=(1,2), \qquad a^2=(1,1),
\]
with prices
\[
q_1=3, \qquad q_2=2.
\]

To find the price of one unit of state-1 consumption, construct the Arrow security $(1,0)$. Observe that
\[
2a^2-a^1 = 2(1,1)-(1,2) = (1,0).
\]
Hence its no-arbitrage price is
\[
2q_2-q_1 = 2\cdot 2 - 3 = 1.
\]
So
\[
p_1=1.
\]

To find the price of one unit of state-2 consumption, construct
\[
a^1-a^2 = (1,2)-(1,1) = (0,1),
\]
whose price is
\[
q_1-q_2 = 3-2 = 1.
\]
So
\[
p_2=1.
\]

Therefore the state price vector is
\[
\boxed{p=(1,1).}
\]
\end{examplebox}

\begin{examplebox}[title={Worked Example: contingent claims from a risky and a risk-free asset}]
Suppose there are two assets:
\begin{itemize}
    \item a risk-free asset with payoff
    \[
    a^f=(1,1)
    \]
    and price
    \[
    q_f=3;
    \]
    \item a risky asset with payoff
    \[
    a^r=(2,1)
    \]
    and price
    \[
    q_r=4.
    \]
\end{itemize}
If wealth is
\[
m=12,
\]
then spending all wealth on the risk-free asset buys
\[
\frac{12}{3}=4
\]
units, which yields
\[
(4,4).
\]
Spending all wealth on the risky asset buys
\[
\frac{12}{4}=3
\]
units, which yields
\[
(6,3).
\]

If a fraction $\alpha$ of wealth is spent on the risky asset and $1-\alpha$ on the risk-free asset, the payoff is
\[
(4,4)+\alpha\bigl((6,3)-(4,4)\bigr)
=
(4,4)+\alpha(2,-1).
\]
So each increase of $2$ units in state-1 payoff costs a reduction of $1$ unit in state-2 payoff. This is exactly the tradeoff implied by state prices
\[
\boxed{p=(1,2).}
\]
\end{examplebox}

\begin{commonmistake}
Do not confuse asset prices with state prices. Asset prices apply to entire payoff vectors; state prices apply to one unit of payoff in a single state. Asset prices are recovered from state prices by
\[
q(a)=p\cdot a.
\]
\end{commonmistake}

\subsection{Diversification and Risk Sharing}

Asset markets allow agents to reallocate consumption across states. This is valuable when people dislike large payoff fluctuations.

\begin{definition}{Diversification}{ch5-diversification}
Diversification means holding a portfolio of assets rather than concentrating wealth in one risky payoff. The aim is to reduce the variance of payoffs across states without sacrificing too much expected return.
\end{definition}

The most effective diversification combines assets with offsetting payoffs, ideally negatively correlated assets.

\begin{examplebox}[title={Worked Example: diversification can create a safe payoff}]
Suppose
\[
a=(1,0),
\qquad
a'=(0,1),
\qquad
\pi=\left(\frac12,\frac12\right).
\]
Each asset alone is risky:
\begin{itemize}
    \item $a$ pays only in state 1;
    \item $a'$ pays only in state 2.
\end{itemize}
But the portfolio
\[
a+a'=(1,1)
\]
is risk-free.

A second example: let
\[
\tilde a=(1,0),
\qquad
\tilde a'=(2,1).
\]
Both are risky, but if shorting is allowed,
\[
\tilde a' - \tilde a = (2,1)-(1,0)=(1,1),
\]
which is again risk-free.

So diversification works not only by buying several assets, but also by combining long and short positions when markets permit it.
\end{examplebox}

\begin{policynote}[title={Policy Application: village risk sharing}]
Suppose each villager faces idiosyncratic crop-failure risk. If one villager's crops fail while others' do not, the villagers can insure one another by trading state-contingent claims. This is risk sharing: each person gives up some consumption in good states in exchange for more consumption in bad states.

However, if \emph{everyone} is hit by the same shock, such as a flood or drought, then internal insurance becomes ineffective. This is \emph{correlated risk}. It explains why some risks, especially large-scale natural-disaster risks, are difficult or expensive to insure.
\end{policynote}

\subsection{Expected Value and Expected Utility}

To rank contingent bundles, the lecture introduces expected value and expected utility.

\begin{definition}{Expected value}{ch5-expected-value}
Given contingent bundle
\[
x=(x_1,\dots,x_S)
\]
with probabilities
\[
\pi=(\pi_1,\dots,\pi_S),
\]
its expected value is
\[
E[x] = \sum_{s=1}^S \pi_s x_s.
\]
\end{definition}

Expected value is a natural benchmark, but by itself it does not capture attitudes toward risk. Many people prefer a safer bundle with lower expected value to a riskier bundle with higher expected value.

\begin{definition}{Expected utility and Bernoulli utility}{ch5-expected-utility}
Let
\[
u:\mathbb{R}_+ \to \mathbb{R}
\]
be a Bernoulli utility function. The expected utility of contingent bundle
\[
x=(x_1,\dots,x_S)
\]
is
\[
U(x)=E[u(x)] = \sum_{s=1}^S \pi_s u(x_s).
\]
\end{definition}

The lecture assumes that $u$ is increasing, and often also concave.

\begin{definition}{Increasing and concave Bernoulli utility}{ch5-inc-concave}
A Bernoulli utility function $u$ is:
\begin{itemize}
    \item \emph{increasing} if more money is better:
    \[
    x>x' \implies u(x)>u(x');
    \]
    \item \emph{concave} if marginal utility is diminishing.
\end{itemize}
If $u$ is differentiable and concave, then $u'$ is decreasing.
\end{definition}

\begin{theorem}{Interior optimality under expected utility}{ch5-focexpected}
Suppose the consumer maximizes
\[
U(x)=\sum_{s=1}^S \pi_s u(x_s)
\]
subject to
\[
p\cdot x \le m,
\]
and the optimum is interior. Then for any two states $i,j$,
\[
\frac{\pi_i u'(x_i)}{\pi_j u'(x_j)} = \frac{p_i}{p_j}.
\]
\end{theorem}

\begin{proof}
The marginal utility of state-$s$ consumption is
\[
MU_s = \pi_s u'(x_s).
\]
At an interior optimum, the tangency condition from earlier chapters gives
\[
\frac{MU_i}{MU_j} = \frac{p_i}{p_j}.
\]
Substituting the marginal utilities yields
\[
\frac{\pi_i u'(x_i)}{\pi_j u'(x_j)} = \frac{p_i}{p_j}.
\]
\end{proof}

This can be rewritten as
\[
\frac{u'(x_i)}{u'(x_j)} = \frac{p_i/\pi_i}{p_j/\pi_j}.
\]
So what matters is the price of state consumption relative to its probability.

\begin{examplebox}[title={Worked Example: optimal contingent consumption with logarithmic utility}]
Suppose there are two states with probabilities
\[
\pi_1=\pi_2=\frac12,
\]
state prices
\[
p=(1,2),
\]
wealth
\[
m=12,
\]
and Bernoulli utility
\[
u(x)=\ln x.
\]
Then expected utility is
\[
U(x_1,x_2)=\frac12 \ln x_1 + \frac12 \ln x_2.
\]

At an interior optimum,
\[
\frac{\pi_1 u'(x_1)}{\pi_2 u'(x_2)}=\frac{p_1}{p_2}.
\]
Since
\[
u'(x)=\frac{1}{x}
\quad \text{and} \quad
\pi_1=\pi_2,
\]
this becomes
\[
\frac{1/x_1}{1/x_2} = \frac{1}{2}
\quad \Longrightarrow \quad
\frac{x_2}{x_1}=\frac12.
\]
So
\[
x_2=\frac{x_1}{2}.
\]
Using the budget constraint
\[
x_1 + 2x_2 = 12,
\]
we obtain
\[
x_1 + 2\left(\frac{x_1}{2}\right)=12
\quad \Longrightarrow \quad
2x_1=12
\quad \Longrightarrow \quad
x_1=6,
\]
and hence
\[
x_2=3.
\]

Therefore the optimal contingent bundle is
\[
\boxed{x=(6,3).}
\]
\end{examplebox}

\subsection{First-Order Stochastic Dominance}

Expected utility with increasing Bernoulli utility respects a natural ranking of risky prospects.

\begin{definition}{Survival probability and first-order stochastic dominance}{ch5-fosd}
Given contingent bundle
\[
x=(x_1,\dots,x_S)
\]
and probability vector $\pi$, define
\[
\pi(x,\alpha)
\]
to be the probability that $x$ pays at least $\alpha$.

We say that $x$ \emph{first-order stochastically dominates} $x'$, written
\[
x \succeq_{FOSD} x',
\]
if
\[
\pi(x,\alpha) \ge \pi(x',\alpha)
\qquad \text{for every } \alpha.
\]
\end{definition}

This means that for every threshold $\alpha$, the probability of receiving at least $\alpha$ is weakly higher under $x$ than under $x'$.

\begin{theorem}{FOSD theorem}{ch5-fosd-theorem}
Suppose preferences take the expected-utility form
\[
U(x)=\sum_{s=1}^S \pi_s u(x_s),
\]
where $u$ is increasing. If
\[
x \succeq_{FOSD} x',
\]
then
\[
U(x)\ge U(x').
\]
\end{theorem}

So FOSD is a robust welfare ranking: every consumer with increasing expected utility agrees that the dominating bundle is at least as good.

\begin{examplebox}[title={Worked Example: checking first-order stochastic dominance}]
Let
\[
\pi=\left(\frac14,\frac14,\frac12\right),
\qquad
x=(7,5,1),
\qquad
x'=(1,0,4).
\]
To test whether
\[
x \succeq_{FOSD} x',
\]
it is enough to check the payoff thresholds appearing in $x'$.

For $\alpha=1$:
\[
\pi(x,1)=1,
\qquad
\pi(x',1)=\frac34.
\]

For $\alpha=0$:
\[
\pi(x,0)=1,
\qquad
\pi(x',0)=1.
\]

For $\alpha=4$:
\[
\pi(x,4)=\frac12,
\qquad
\pi(x',4)=\frac12.
\]

Thus
\[
\pi(x,\alpha)\ge \pi(x',\alpha)
\]
for all relevant thresholds, so
\[
\boxed{x \succeq_{FOSD} x'.}
\]
Therefore every consumer with increasing expected utility weakly prefers $x$ to $x'$.
\end{examplebox}

\subsection{Risk Attitudes}

Concavity of Bernoulli utility captures dislike of risk.

\begin{theorem}{Jensen inequality for expected utility}{ch5-jensen}
If $u$ is concave, then for every contingent bundle $x$,
\[
u(E[x]) \ge E[u(x)].
\]
Equivalently,
\[
u(E[x]) \ge U(x).
\]
\end{theorem}

\begin{proof}
This is Jensen's inequality applied to the concave function $u$.
\end{proof}

The economic meaning is crucial: if utility is concave, then receiving the expected value for sure is at least as good as facing the risky prospect itself.

\begin{definition}{Risk attitudes}{ch5-risk-attitudes}
Let
\[
U(x)=E[u(x)].
\]
Then:
\begin{itemize}
    \item the consumer is \emph{risk averse} if
    \[
    u(E[x]) \ge U(x)
    \quad \text{for all } x;
    \]
    \item the consumer is \emph{risk neutral} if
    \[
    u(E[x]) = U(x)
    \quad \text{for all } x;
    \]
    \item the consumer is \emph{risk loving} if
    \[
    u(E[x]) \le U(x)
    \quad \text{for all } x.
    \]
\end{itemize}
\end{definition}

\begin{theorem}{Linear utility implies risk neutrality}{ch5-risk-neutral}
A consumer is risk neutral if the Bernoulli utility function is affine:
\[
u(x)=ax+b
\qquad \text{with } a>0.
\]
In particular, $u(x)=x$ is risk neutral.
\end{theorem}

\begin{proof}
If
\[
u(x)=ax+b,
\]
then
\[
U(x)=E[u(x)] = E[ax+b] = aE[x]+b = u(E[x]).
\]
So the consumer is risk neutral.
\end{proof}

\begin{examplebox}[title={Worked Example: same expected value, different risk}]
Suppose
\[
\pi=\left(\frac12,\frac12\right).
\]
Compare the safe bundle
\[
x=(10^7,10^7)
\]
with the risky bundle
\[
x'=(4\times 10^7,0).
\]
Their expected values are
\[
E[x]=10^7,
\qquad
E[x']=2\times 10^7.
\]
So $x'$ has the higher expected value.

However, if the consumer has increasing and concave Bernoulli utility, expected utility need not follow expected value. A sufficiently risk-averse consumer may prefer the certain bundle $x$ to the much riskier $x'$ despite its lower expected value.

This is exactly why expected utility rather than expected value is the central decision criterion in the lecture.
\end{examplebox}

\begin{examplebox}[title={Worked Example: certainty equivalent under concavity}]
Suppose
\[
\pi=\left(\frac12,\frac14,\frac14\right),
\qquad
x=(4,8,12).
\]
Then the expected value is
\[
E[x]=\frac12\cdot 4 + \frac14\cdot 8 + \frac14\cdot 12 = 2+2+3=7.
\]
If $u$ is increasing and concave, Jensen's inequality gives
\[
u(7) = u(E[x]) \ge U(x).
\]
So the consumer weakly prefers receiving $7$ for sure to receiving the contingent bundle $(4,8,12)$ itself.

This is the standard exam interpretation of concavity:
\[
\boxed{\text{concavity } \Longleftrightarrow \text{ risk aversion}.}
\]
\end{examplebox}

\begin{commonmistake}
Do not conclude that a bundle with higher expected value must be preferred. That is true only for risk-neutral consumers. Under concave expected utility, a lower-mean but less risky bundle can be preferred.
\end{commonmistake}

\subsection{Chapter Summary}

\begin{keyformula}[title={Key Formula: Lecture 5 Summary}]
\begin{align*}
x &= (x_1,\dots,x_S) && \text{contingent consumption bundle}, \\
p\cdot x &\le m && \text{state-contingent budget}, \\
q\cdot \theta &= \sum_{k=1}^K q_k \theta_k && \text{portfolio cost}, \\
\sum_{k=1}^K \theta_k a^k &&& \text{portfolio payoff}, \\
q(\tilde a) &= q\cdot \theta && \text{if } \tilde a=\sum_{k=1}^K \theta_k a^k, \\
q(a) &= p\cdot a && \text{state-price representation}, \\
E[x] &= \sum_{s=1}^S \pi_s x_s, \\
U(x) &= E[u(x)] = \sum_{s=1}^S \pi_s u(x_s), \\
\frac{\pi_i u'(x_i)}{\pi_j u'(x_j)} &= \frac{p_i}{p_j} && \text{interior optimum}, \\
x \succeq_{FOSD} x' &\Longrightarrow U(x)\ge U(x') && \text{if } u \text{ is increasing}, \\
u(E[x]) &\ge E[u(x)] && \text{if } u \text{ is concave}.
\end{align*}
\end{keyformula}

\begin{examtip}
The highest-yield skills from Lecture 5 are:
\begin{itemize}
    \item represent uncertain outcomes as contingent consumption bundles;
    \item translate insurance or asset trading into state-contingent budgets;
    \item identify arbitrage and use no-arbitrage to price replicated assets;
    \item recover state prices from asset prices in simple two-state settings;
    \item use expected utility rather than expected value when risk attitudes matter;
    \item test first-order stochastic dominance carefully using payoff thresholds;
    \item interpret concavity of Bernoulli utility as risk aversion.
\end{itemize}
\end{examtip}